% Options for packages loaded elsewhere
\PassOptionsToPackage{unicode}{hyperref}
\PassOptionsToPackage{hyphens}{url}
%
\documentclass[
]{article}
\usepackage{amsmath,amssymb}
\usepackage{lmodern}
\usepackage{iftex}
\ifPDFTeX
  \usepackage[T1]{fontenc}
  \usepackage[utf8]{inputenc}
  \usepackage{textcomp} % provide euro and other symbols
\else % if luatex or xetex
  \usepackage{unicode-math}
  \defaultfontfeatures{Scale=MatchLowercase}
  \defaultfontfeatures[\rmfamily]{Ligatures=TeX,Scale=1}
\fi
% Use upquote if available, for straight quotes in verbatim environments
\IfFileExists{upquote.sty}{\usepackage{upquote}}{}
\IfFileExists{microtype.sty}{% use microtype if available
  \usepackage[]{microtype}
  \UseMicrotypeSet[protrusion]{basicmath} % disable protrusion for tt fonts
}{}
\makeatletter
\@ifundefined{KOMAClassName}{% if non-KOMA class
  \IfFileExists{parskip.sty}{%
    \usepackage{parskip}
  }{% else
    \setlength{\parindent}{0pt}
    \setlength{\parskip}{6pt plus 2pt minus 1pt}}
}{% if KOMA class
  \KOMAoptions{parskip=half}}
\makeatother
\usepackage{xcolor}
\setlength{\emergencystretch}{3em} % prevent overfull lines
\providecommand{\tightlist}{%
  \setlength{\itemsep}{0pt}\setlength{\parskip}{0pt}}
\setcounter{secnumdepth}{-\maxdimen} % remove section numbering
\ifLuaTeX
  \usepackage{selnolig}  % disable illegal ligatures
\fi
\IfFileExists{bookmark.sty}{\usepackage{bookmark}}{\usepackage{hyperref}}
\IfFileExists{xurl.sty}{\usepackage{xurl}}{} % add URL line breaks if available
\urlstyle{same} % disable monospaced font for URLs
\hypersetup{
  hidelinks,
  pdfcreator={LaTeX via pandoc}}

\author{Dietmar Gerald Schrausser}
\date{09.05.2026}


\begin{document}

\hypertarget{foliumblat-vv}{%
\section{Foliu{[}m{]}/Blat Vv}\label{foliumblat-vv}}

A \emph{contemporary} English translation is presented from comparing
the existing Latin, Early New High German (ENHG) and English
(translation) texts for better comprehensibility. Especially since the
ENHG text is difficult to understand compared to modern High German, as
a \emph{profound} knowledge of German (as a mother tongue) as well as
already aknowledge of various \emph{dialects} is required.

\hypertarget{of-the-sanctification-of-the-seventh-day}{%
\subsection{Of the Sanctification of the Seventh
Day}\label{of-the-sanctification-of-the-seventh-day}}

When the world, a work of \emph{divine wisdom}, was completed in six
days, heaven and earth created, ordered, and finally adorned, the
glorious God finished his work and rested from his (manual) labors on
the \emph{seventh} day. This completion, however, did not mean an
\emph{ending}, as if he were \emph{tired} of the work, but rather to
\emph{not}\footnote{`facere cessauit' and `zemachen {[}\ldots{]} nit
  vergange{[}n{]}'.} create any further creature of other matter or
likenesses, especially since the work of creation does not cease to have
an effect.\footnote{from Petrus Lombardus' `Sententiarum' (Petrus,
  \href{https://doi.org/10.3931/e-rara-18241}{1484}, Book II, Chapter7).}

The Lord blessed and sanctified \emph{this} day and called it
\emph{Sabbatum}, which in Hebrew means \emph{rest}\footnote{`requiem'
  and `ein rüe'.}, as he rested on that day from all his work.

\begin{quote}
Eo q{[}uia{]} in ipso cessauerat ab om{[}n{]}i opere q{[}uo{]}d
patrarat.\footnote{`For on this day he had ceased all his works,'
  `Vulgate', a variation of Genesis 2:3 (c.f. Jiménez de Cisneros,
  \href{https://doi.org/10.3931/e-rara-46695}{1517}).}
\end{quote}

It is also customary among Jews to celebrate this day by abstaining from
work and it was also \emph{law} among many pagan peoples to celebrate
this day.

Thus we have reached the end of the \emph{divine works}. Accordingly, we
should (i) \emph{approach with loving reverence}\footnote{perhaps better
  than `fear, love and honor'.}, where all things visible and invisible
exist, (ii) seek\footnote{`queramus' and `suchen'.} present
\emph{values}\footnote{`presentia bona: quatenus bona sint' and `die
  gegenwürtigen güter. sover die gut sind', refers to \emph{present
  blessings}, in a philosophical or theological context, present,
  material or earthly things, as opposed to eternal, spiritual goods.
  (c.f. Thomas Aquinas,
  \href{https://doi.org/10.3931/e-rara-15170}{1485} or Augustine,
  \href{https://archive.org/details/expositionsonboo39augu/mode/1up}{1847}).}
(insofar as \emph{good}\footnote{\emph{benevolent, humanitarian,
  philanthropic, gracious, charitable}.}) and the true
\emph{blessedness} of eternal life from the Lord of Heaven, the Lord of
all good things, to whom power in heaven and on earth has been given.

\hypertarget{references}{%
\subsection{References}\label{references}}

Augustinus, A. (1847). \emph{Expositions on the Book of Psalms}. Oxford:
J. H. Parker.
\url{https://archive.org/details/expositionsonboo39augu/mode/1up}

Jiménez de Cisneros, F. (1517). \emph{Biblia Polyglotta Complutensis}.
Complutum: Arnaldo Guillén de Brocar.
\url{https://doi.org/10.3931/e-rara-46695}

Petrus. (1484). \emph{Sententiarum Libri IV}. Basel.
\url{https://doi.org/10.3931/e-rara-18241}

Thomas. (1485). \emph{Summa Theologica: Partes 1-3}. Basel: Michael
Wenssler. \url{https://doi.org/10.3931/e-rara-15170}

\end{document}
